%% GENERIC STYLE SETTINGS
\usetheme{default}
\usepackage{listings}
\usepackage{multirow}
\usepackage{xcolor}

\usebackgroundtemplate{
\includegraphics[width=\paperwidth,height=\paperheight]{images/hutton_background}
}
%% PRESENTATION CONFIGURATION PARAMETERS %%%%%%%%%%%%%%%%%%%%%%%%%%%%%%%%%%%%%%%
%\titlebackgroundfile{images/hutton_title}
%\framebackgroundfile{images/hutton_background}
\definecolor{hutton_green}{HTML}{78A22F}
\definecolor{hutton_purple}{HTML}{872175}
\definecolor{hutton_blue}{HTML}{569BBE}
\usefonttheme{structurebold}
\setbeamercolor{alerted text}{fg=orange}
\setbeamercolor{background canvas}{bg=white}
\setbeamercolor{block title}{bg=hutton_purple}
\setbeamercolor{frametitle}{fg=hutton_purple}
\setbeamercolor{title}{fg=black}
\setbeamercolor{titlelike}{fg=hutton_green}
\setbeamercolor{author}{fg=hutton_purple}
\setbeamercolor{author in head/foot}{fg=white}
\setbeamercolor{title in head/foot}{fg=white}
\setbeamercolor{section in head/foot}{fg=hutton_purple}
\setbeamercolor{normal text}{fg=black}
\setbeamercolor{frametitle}{fg=hutton_purple}
\setbeamerfont{block title}{size={}}
\setbeamerfont{author}{size=\footnotesize}
\setbeamerfont{date}{size=\footnotesize}
\setbeamercolor{section in toc shaded}{fg=hutton_purple}
\setbeamercolor{section in toc}{fg=hutton_purple}
\setbeamercolor{subsection in toc shaded}{fg=hutton_purple}
\setbeamercolor{subsection in toc}{fg=hutton_purple}
\setbeamertemplate{itemize item}[circle]
\setbeamertemplate{itemize subitem}[circle]
\setbeamertemplate{itemize subsubitem}[circle]
\setbeamertemplate{itemize subsubsubitem}[circle]
\setbeamercolor{itemize item}{fg=hutton_purple}
\setbeamercolor{itemize subitem}{fg=hutton_purple}
\setbeamercolor{itemize subsubitem}{fg=hutton_purple}
\setbeamercolor{itemize subsubsubitem}{fg=hutton_purple}
\setbeamercolor{enumerate item}{fg=hutton_purple}
\setbeamercolor{enumerate subitem}{fg=hutton_purple}
\setbeamercolor{enumerate subsubitem}{fg=hutton_purple}
\setbeamercolor{enumerate subsubsubitem}{fg=hutton_purple}
\setbeamercolor{alerted text}{fg=hutton_green}
\setbeamerfont{alerted text}{series=\bfseries}
% This command makes sure that acrobat reader doesn't change the colours of the slide
% when there are figures with transparencies.
\pdfpageattr {/Group << /S /Transparency /I true /CS /DeviceRGB>>}

%Disables discrete bottom navigation bar
%\beamertemplatenavigationsymbolsempty

% Modify the slide titles to avoid the corner images,
\setbeamertemplate{frametitle}
{
\vspace{0.05\textheight}
\noindent\quad\begin{minipage}[t][0.12\textheight][t]{0.85\textwidth}
\insertframetitle\par
\end{minipage}
}

% Modify title page to avoid the big logo on right
\setbeamertemplate{title page}{
    \begin{picture}(0,0)
            %This ends up on top of the default background image, rather than replacing it:
            \put(-30,-165){%
                \includegraphics[width=\paperwidth,height=\paperheight]{images/hutton_title}
            }
            \put(0,-75){%
                \begin{minipage}[b][0.4\textheight][t]{0.75\textwidth}
                    \usebeamerfont{title}\usebeamercolor[fg]{title}{\inserttitle\par}
                    \usebeamerfont{subtitle}\usebeamercolor[fg]{subtitle}{\insertsubtitle\par}
                \end{minipage}
            }
            \put(0,-135){%
                \begin{minipage}[b][0.1\textheight][t]{\textwidth}
                    \usebeamerfont{author}\usebeamercolor[fg]{author}{\insertauthor\par}
                \end{minipage}
            }
    \end{picture}
}

%%%%%%%%%%%%%%%%%%%%%%%%%%%%%%%%%%%%%%%%%%%%%%%%%%%%%%%%%%%%%%%%%%%%%%%%%%%%%%%%

% LISTINGS SETTING
% Settings for code listings in lstlistings

\definecolor{hutton_lightgreen}{HTML}{C8F27F}

\lstset{ %
  backgroundcolor=\color{hutton_lightgreen},   % choose the background color; you must add \usepackage{color} or \usepackage{xcolor}
  basicstyle=\tiny\ttfamily,        % the size of the fonts that are used for the code
  breakatwhitespace=true,          % sets if automatic breaks should only happen at whitespace
  breaklines=true,                 % sets automatic line breaking
  captionpos=b,                    % sets the caption-position to bottom
  commentstyle=\color{red},    % comment style
  deletekeywords={...},            % if you want to delete keywords from the given language
  escapeinside={\%*}{*)},          % if you want to add LaTeX within your code
  extendedchars=true,              % lets you use non-ASCII characters; for 8-bits encodings only, does not work with UTF-8
  frame=single,                    % adds a frame around the code
  keepspaces=true,                 % keeps spaces in text, useful for keeping indentation of code (possibly needs columns=flexible)
  keywordstyle=\color{blue},       % keyword style
%  language=Octave,                 % the language of the code
  morekeywords={*,...},            % if you want to add more keywords to the set
  numbers=left,                    % where to put the line-numbers; possible values are (none, left, right)
  numbersep=5pt,                   % how far the line-numbers are from the code
  numberstyle=\tiny\color{gray}, % the style that is used for the line-numbers
  rulecolor=\color{black},         % if not set, the frame-color may be changed on line-breaks within not-black text (e.g. comments (green here))
  showspaces=false,                % show spaces everywhere adding particular underscores; it overrides 'showstringspaces'
  showstringspaces=false,          % underline spaces within strings only
  showtabs=false,                  % show tabs within strings adding particular underscores
  stepnumber=1,                    % the step between two line-numbers. If it's 1, each line will be numbered
  stringstyle=\color{violet},     % string literal style
  tabsize=4,                       % sets default tabsize to 2 spaces
  title=\lstname                   % show the filename of files included with \lstinputlisting; also try caption instead of title
}


%%%
% TITLE PREAMBLE
\title[Open science and Bio* projects] % (optional, only for long titles)
{Open science and Bio* projects}
\subtitle{Standing on each other's shoulders}
\author[Cock] % (optional, for multiple authors)
{Peter~Cock (\href{https://twitter.com/pjacock}{@pjacock on Twitter})}
\institute[The James Hutton Institute] % (optional)
{
  Information and Computational Sciences\\
  The James Hutton Institute
}
\date[July 2017] % (optional)
{EASTBIO Genomic Approaches, 1$^{st}$ July 2017}
\subject{Bioinformatics}

%%%
% TOC
% Show table of contents, with current section highlighted,
% at the start of each section

\AtBeginSection[]
{
  \begin{frame}
    \frametitle{Table of Contents}
    \tableofcontents[currentsection,hideallsubsections]
  \end{frame}
}

%%%
% START DOCUMENT
\begin{document}

\frame[plain]{\titlepage}

\section{Introduction}
%\subsection{BLAST Extension}

\begin{frame}
  \frametitle{Progress in Science}
  \begin{itemize}
    \item \emph{``If I have seen further it is by standing on the shoulders of Giants.''} (Issac Newton, 1676)
    \item Science works best when can build on past work
    \item Methods section in papers \emph{should} let you do this
    \item For computation methods, would ideally get the code...
  \end{itemize}
\end{frame}

\begin{frame}
  \frametitle{Progress in Open Source}
  \begin{itemize}
    \item By default, computer code is under copyright:
      \begin{itemize}
        \item YOU NEED PERMISSION to use, re-use, adapt.
      \end{itemize}
    \item Open Source licensing defines explicit terms:
      \begin{itemize}
        \item GIVES PERMISSION to use, re-use, share, change, at no cost.
        \item That includes packaging for easier installation
      \end{itemize}
  \end{itemize}
\end{frame}

\section{Open Source}

\subsection{Open Source Licenses}

\begin{frame}
  \frametitle{Many Open Source Licenses}
  \begin{itemize}
    \item Open Source Initiative \url{https://opensource.org/} has a list
    \item Very liberal, e.g. MIT and BSD
    \item `Viral', e.g. GNU Public Licence (GPL)
  \end{itemize}
  Note: Combinging code from multiple  sources can hit license problems
\end{frame}

\subsection{Code sharing}

\begin{frame}
  \frametitle{Code sharing / development}
  Most projects use a web-based platform:
  \begin{itemize}
    \item \href{https://bitbucket.org/product}{Bitbucket}
    \item \href{https://github.com/}{GitHub}
    \item \href{https://about.gitlab.com/}{GitLab}
    \item Sourceforge (popular over 10 years ago)
    \item Lab specific website
  \end{itemize}
\end{frame}

\begin{frame}
  \frametitle{GitHub/BitBucket's Business Model}
  \begin{itemize}
    \item Offer their services for free to \emph{public} projects
    \item Offer private projects for free to \emph{students} \& \emph{teachers}
    \item Charge \$\$\$ for private projects
  \end{itemize}
  i.e. We're freeloading and can't count of this continuing.

  Compare with SourceForge, which is now covered in adverts
\end{frame}

\subsection{Open Source in Biology}

\begin{frame}
  \frametitle{Open Source in Biology}
  \begin{itemize}
    \item Much early bioinformatics by Systems Administrators
    \item e.g. \emph{How Perl Saved the Human Genome Project} (Lincoln Stein, 1996)
    \item Absorbed ideas/culture from free and open source world, e.g. Linux, Perl
    \item Open source required championing in early days, e.g. \href{https://www.open-bio.org/wiki/BOSC}{Bioinformatics Open Source Conference} since 2000.\footnote{Peter is a former co-chair}
    \item Many journals now on-board, e.g.
    \href{http://journals.plos.org/plosone/s/materials-and-software-sharing}{PLOS},
    \href{https://academic.oup.com/nar/pages/Policies}{NAR},
    \href{https://academic.oup.com/gigascience/pages/editorial_policies_and_reporting_standards}{GigaScience}
  \end{itemize}
\end{frame}

\begin{frame}
  \frametitle{Open Source Bioinformatics examples}
  Important funded projects in genomics:
  \begin{itemize}
    \item \url{https://github.com/samtools/samtools} \\ Wellcome Trust Sanger Institute
    \item \url{https://github.com/broadinstitute/gatk} \\ Broad Institute
    \item \url{https://github.com/galaxyproject/galaxy} \\ Penn State
  \end{itemize}
\end{frame}

\begin{frame}
  \frametitle{Open Source Bioinformatics examples}
  Large long-running un-funded projects:
  \begin{itemize}
    \item \url{https://github.com/bioperl/bioperl-live} - \href{http://bioperl.org}{BioPerl}
    \item \url{https://github.com/biopython/biopython} - \href{http://biopython.org}{Biopython}\footnote{Peter is a major contributor}
    \item \url{https://github.com/bioruby/bioruby} - \href{http://bioruby.org}{BioRuby}
    \item \url{https://github.com/biojava/biojava} - \href{http://biojava.org}{BioJava}
  \end{itemize}
\end{frame}

\begin{frame}
  \frametitle{Open Source Bioinformatics examples}
  Relevant smaller projects, by individuals or smaller groups:
  \begin{itemize}
    \item \url{https://github.com/arq5x/bedtools}
    \item \url{https://github.com/arq5x/poretools}
    \item \url{https://sourceforge.net/projects/mira-assembler/}
    \item \url{https://github.com/tseemann/prokka}
    \item \url{https://github.com/COMBINE-lab/salmon}
  \end{itemize}
\end{frame}

\begin{frame}
  \frametitle{Open Source Bioinformatics examples}
  Sequencing companies with open source projects:
  \begin{itemize}
    \item \url{https://github.com/illumina} \\ Illumina
    \item \url{https://github.com/nanoporetech/} \\ Oxford Nanopore Technologies
    \item \url{https://github.com/PacificBiosciences} \\ Pacific Biosciences
    \item \url{https://github.com/iontorrent} \\ Ion Torrent
  \end{itemize}
\end{frame}

\begin{frame}
  \frametitle{GitHub is currently very popular}
  \begin{itemize}
    \item Has built up a network effect, the place to go
    \item Continues to innovate and improve the platform
    \item Hub of an ecosystem with other services (e.g. testing)
    \item Also being used for teaching materials, slides, etc
    \item Doubles as an online profile for individuals
    \item Thus far, free for open projects...
  \end{itemize}
\end{frame}

\begin{frame}
  \frametitle{Publishing Source code for long term?}
  \begin{itemize}
    \item Cannot count on GitHub/Bitbucket for long term storage
    \item Personal webpages even more fragile
    \item Scientific record needs long term static archive:
      \begin{itemize}
        \item Journal paper supplementary material
        \item \href{http://about.zenodo.org/}{Zenodo} (powered by CERN)
      \end{itemize}
  \end{itemize}
\end{frame}

\section{Open Data}

\begin{frame}
  \frametitle{Open Data in Genomics}
  Sequence data is \emph{usually} shared, now expected:
  \begin{itemize}
    \item DNA, RNA, Proteins at EMBL-EBI, NCBI, DDBJ
    \item Raw sequencing data too at EMBL SRA, NCBI SRA
    \item Structures at RSCB Protein Data Bank
  \end{itemize}
  Most journals will enforce this.
\end{frame}

\begin{frame}
  \frametitle{Open Data in Genomics}
  Curated data is not always freely openly shared,
\begin{columns}[T]
\begin{column}{.5\textwidth}
%\begin{block}{Image}
  \begin{itemize}
    \item Kyoto Encyclopedia of Genes and Genomes (KEGG) now has \href{http://www.kegg.jp/kegg/legal.html}{restrictions}
    \item The Arabidopsis Information Resource (TAIR) now has \href{https://www.arabidopsis.org/doc/about/tair_subscriptions/413}{fees}
    \item Also some groups like to hoard their precious data...
  \end{itemize}
%\end{block}
\end{column}
\begin{column}{.5\textwidth}
%\begin{block}{Text}
  \includegraphics[width=\textwidth]{images/gollum.jpg}
%\end{block}
\end{column}
\end{columns}
\end{frame}

\section{Open Data \& Software $\rightarrow$ Reproducible Science?}

\begin{frame}
  \frametitle{Open Data \& Software $\rightarrow$ Reproducible?}
  Given a well written scientific paper:
  \begin{itemize}
    \item The raw data will be available (e.g. NCBI accessions)
    \item The tools used and exact versions will be stated
    \item The actual workflow, scripts or command lines are given
  \end{itemize}
  How easy would it be to repeat the analysis?
\end{frame}


\begin{frame}
  \frametitle{Open Data \& Software $\rightarrow$ Reproducible?}
  \begin{itemize}
    \item Sadly click-and-repeat methods in papers are still an aspiration
    \item Being worked on, still lots of technical hurdles
      \begin{itemize}
        \item Including installing dependencies, but see e.g. BioConda
      \end{itemize}
    \item However, openly sharing the data and code is first step!
  \end{itemize}
\end{frame}

\section{Conclusions}

\subsection{Sources of advice}

\begin{frame}
  \frametitle{Sources of coding / tool advice}
  \begin{itemize}
    \item Ask us your peers / supervisor in person
    \item For a given tool check its documentation, mailing lists, issue tracker, online chat/forum, etc
    \item BioStars \url{https://www.biostars.org}
    \item Bioinformatics StackExchange \url{https://bioinformatics.stackexchange.com/}
    \item SeqAnswers Forums \url{http://seqanswers.com/forums}
 \end{itemize}
\end{frame}

\begin{frame}
  \frametitle{Conclusions}
  \begin{itemize}
    \item Open science is a community - help each other
    \item Register a GitHub account to contribute back
      \begin{itemize}
        \item Raise quries; log bugs; offer improvements
        \item Use this to share your own scripts/tools
      \end{itemize}
    \item Your GitHub (etc) profile can be part of your CV, e.g.
      \begin{itemize}
        \item \url{https://github.com/peterjc} - my GitHub profile
        \item \url{https://bitbucket.org/peterjc/} - my BitBucket profile
        \item \url{https://gitlab.com/pjacock} - my GitLab profile
        \item \url{https://www.biostars.org/u/146/} - my BioStars profile
      \end{itemize}
 \end{itemize}
\end{frame}

\subsection{Further reading}

\begin{frame}
  \frametitle{Further reading}
  \begin{itemize}

    \item Mike Croucher \emph{Is your research software correct?}
    \url{https://mikecroucher.github.io/MLPM_talk/}
    % inc. Croucher's Law \emph{I can be an idiot and will make mistakes}

    \item Software Sustainability Institute
    \url{https://www.software.ac.uk/}

    \item Software Carpentry
    \url{https://software-carpentry.org/}

    \item Wilson et al (2014) Best Practices for Scientific Computing
    \url{https://doi.org/10.1371/journal.pbio.1001745}

  \end{itemize}
\end{frame}

% etc
\end{document}
